Запишем важную и полезную функцию, которая позволяет извлекать i-й элемент маски:
\begin{lstlisting}
bool bit(long long mask, int pos) {
  return (mask >> pos) & 1;
}
\end{lstlisting}
\textbf{Формулировка задачи} \\ Дан полный, неориентированный, взвешенный граф.\\
\textit{Гамильтонов путь в графе} - это путь, который начинается в какой-то вершине графа, проходит по всем вершинам и посещает все вершины графа ровно по одному разу

Среди всех таких путей нас интересует путь минимальной стоимости

Пусть у нас есть какой-то путь, проходящий по каким-то вершинам графа ровно один раз и заканчивающийся в вершинке v. Тогда для того, чтобы продолжить этот путь, нам необходимо знать, в какой вершинке мы находимся и маску посещенный вершин. Отсюда возникает dp:
\begin{enumerate}
    \item $dp[v][mask]$ - это стоимость минимального пути, который где-то начинается, посещает все вершины маски ровно по одному разу и заканчивается в $v$
    \item База: Пути длины 0. Когда мы стоим в какой-то вершинке и еще никуда не пошли. \\ $dp[v][2^v]\,=\,0$, все остальные значения положим $+\infty$
    \item Переход: перебираем все ребра, исходящие из $v$, которые ведут в вершины, еще не посещенные (в маске на месте соответствующей вершины будет стоять 0), и смотрим, куда можем пойти.
    \item Ответ: Минимальное $dp[u][2^n - 1]$. То есть нас интересуют все пути, которые посещают все вершины, выбираем из них тот, что имеет минимальную стоимость
\end{enumerate}

 $mask\,|\,(1<<u)$ - взять побитовое "или" числа $2^u$ и $mask$. Нетрудно заметить, что такая процедура просто ставит единичку на u-й бит маски
\begin{lstlisting}
for (int mask = 0; mask < (1 << n); ++mask) {
  for (int v = 0; v < n; ++v) {
    for (int u = 0; u < n; ++u) {
      if (bit(mask,u)) {
        continue;
      }
      int newmask = mask | (1 << u);
      dp[u][newmask] = min(dp[u][newmask], dp[v][mask] + cost[v][u]);
    }
  }
}
\end{lstlisting}

\textbf{Асимптотика} \\
За счет циклов $O(2^n n^2)$

\pagebreak